\cleardoublepage
\phantomsection
\addcontentsline{toc}{part}{Calabi--Yau Quantum Groups}
\chapter*{Book entrance: Calabi--Yau Quantum Groups}
\addcontentsline{toc}{chapter}{Book entrance: Calabi--Yau Quantum Groups}
\begin{summarybox}
An extension of geometric objects defines a correspondence. Its action on cohomology depends on the coefficient theory and the allowed pushforward. Tensor products and exchange require a coproduct and compatible intertwiners. A trace on a specified representation then produces an amplitude. The reconstruction problem is to retain these operations through the comparisons between Calabi--Yau categories, boundary fields, Hall algebras, and quantum groups.
\end{summarybox}
\part{Categorical and field-theoretic foundations}

\chapter{Calabi--Yau categories}
\section{Smoothness, properness, and the Calabi--Yau class}
Let $\Ccal$ be a small stable $\kfield$-linear category. Smoothness means that the diagonal bimodule is perfect. Properness means that the morphism complexes are perfect over $\kfield$. A $d$-Calabi--Yau structure is a nondegenerate negative-cyclic class
\[
 \theta_{\CY}\in HC^-_d(\Ccal)
\]
whose underlying Hochschild trace identifies the diagonal bimodule with its dual shifted by $d$.

For $E\in\Ccal$, the tangent complex of the derived moduli stack of objects is
\[
 T_E\M_{\Ccal}\simeq\RHom_{\Ccal}(E,E)[1].
\]
The Calabi--Yau trace pairs two tangent vectors by composition and trace.

\begin{theorem}[The two universal categorical outputs]
Let $\Ccal$ be a smooth proper $d$-Calabi--Yau category.
\begin{enumerate}
\item The Hochschild cochains $CC^*(\Ccal)$ carry a framed $\Etwo$-algebra structure.
\item Whenever the moduli stack of objects is representable in the PTVV setting, $\M_{\Ccal}$ carries a $(2-d)$-shifted symplectic form.
\end{enumerate}
\end{theorem}
\begin{proof}
The cyclic Deligne theorem upgrades the brace $\Etwo$ structure on Hochschild cochains by the circle action supplied by the cyclic trace \cite{BravRozenblyum}. The shifted form at $E$ is
\[
 \RHom(E,E)[1]^{\tensor2}
 \xrightarrow{\ \circ\ }
 \RHom(E,E)[2]
 \xrightarrow{\ \Tr\ }
 \kfield[2-d].
\]
Nondegeneracy is the Calabi--Yau duality isomorphism, and closedness is the negative-cyclic lift \cite{PTVV,BravDyckerhoff}.
\end{proof}

The integer $d$ controls the shift of the moduli form. The Hochschild operation remains framed $\Etwo$. These structures meet through deformation theory while retaining their separate operadic degrees.

\chapter{Local BV theories and factorization observables}
\section{Cyclic elliptic local Lie algebras}
Let $X$ be a complex manifold. A classical local field theory is specified by a sheaf $\Lcal$ of elliptic $L_\infty$ algebras, a local invariant pairing, and an action functional $S$ satisfying the classical master equation. For an open set $U\subset X$, the classical observables are
\[
 \Obs^{\mathrm{cl}}(U)=C^*_{\CE,\mathrm{cont}}(\Lcal(U)).
\]
Here $\Lcal(U)$ consists of unrestricted smooth fields with its
nuclear smooth-section topology. Continuous cochains use the completed
symmetric powers of its continuous dual, with the completed tensor
products of the elliptic-observable construction \cite{CG1,CG2}.
The dual consists of compactly supported distributions.
For $U\subset V$, restriction of fields $\Lcal(V)\to\Lcal(U)$
induces the covariant map on observables.
For pairwise disjoint $U_i\subset V$, restrict to
$\bigoplus_i\Lcal(U_i)$ and multiply the pulled-back cochains.
These are the prefactorization maps.

\begin{theorem}[Factorization observables]
An elliptic classical BV theory produces a factorization algebra $\Obs^{\mathrm{cl}}$. A renormalized quantization satisfying the quantum master equation produces a factorization algebra $\Obs^{\mathrm{q}}$ over $\kfield[[\hbar]]$.
\end{theorem}
\begin{proof}
Restriction is a morphism of local $L_\infty$ algebras, so its
contravariant continuous Chevalley map commutes with the differential.
For disjoint opens, the completed cochains of the finite direct sum
are the completed tensor product of the cochains. This constructs
the displayed maps and their composition identities.
The elliptic descent theorem in this completed coefficient category
supplies Weiss descent \cite{CG1,CG2}. In the quantum theory, renormalized time-ordered products define the structure maps, and the quantum master equation makes their differential square-zero and compatible with nested ultraviolet contractions \cite{CG1,CG2}.
\end{proof}

\section{Boundary and defect restriction}
Let $Y=C\times\R_{\ge0}$, with $C$ a complex curve, and let a boundary condition be given at $C\times\{0\}$. Boundary observables are holomorphic along $C$ and ordered along the inward normal line. A compact boundary object $b$ in the boundary factorization category gives
\[
 A_b=\RHom(b,b).
\]

\begin{theorem}[Ordered boundary algebra]
When boundary observables are holomorphic on $C$ and locally constant along the oriented normal direction, $A_b$ is an ordered chiral algebra
\[
 A_b\in\Alg_{\Eone}\bigl(\Fact_D(C),\tensor^!\bigr).
\]
If $b$ is a local compact generator and the Morita equivalences obey factorization descent, the boundary category is equivalent to $A_b$-modules in factorization coefficients.
\end{theorem}
\begin{proof}
Composition of boundary endomorphisms gives the internal $\Eone$ multiplication. Holomorphic collision gives the factorization coefficient on $C$. The componentwise locality identity makes multiplication compatible with every disjoint-support factorization map. Local compact-generator Morita theory descends over the Ran diagrams. These are the recognition and Morita theorems of Volume~I.
\end{proof}


\chapter{Three curve-level algebraic outputs}
Let $C$ be a smooth curve and let $\Dcal_C$ be the fixed Ran category.  The same geometric source can produce three algebraic outputs, and their distinction is functorial.

\begin{definition}[Curve-level faces]
A transverse face is
\[
 A^\perp\in\Alg_{\Eone}(\Fact_D(C),\tensor^!).
\]
An associative chiral face is
\[
 A^{\ch}_{\mathsf{Ass}}\in\Alg_{\Eone}(\Dcal_C,\starCh).
\]
A chiral Lie face is
\[
 \g^{\ch}\in\Alg_{\mathsf{Lie}}(\Dcal_C,\starCh).
\]
Their bars are $\BarPerp$, $\BarAssCh$, and $\BarLieCh$.
\end{definition}

\begin{theorem}[Source separation]
A boundary normal direction determines the transverse face.  A chiral-convolution multiplication determines the associative chiral face.  A Lie bracket for chiral convolution determines the Lie face.  The commutator functor sends the associative chiral face to the Lie face.  A chiral enveloping algebra $U^{\ch}(\g)$ satisfies
\[
 \BarAssCh(U^{\ch}(\g))\simeq\BarLieCh(\g)
\]
on the nilcomplete PBW locus.  A mixed source joins the transverse and associative chiral faces through an aligned distributive law.
\end{theorem}
\begin{proof}
Each statement is the corresponding operadic restriction or adjunction in the named monoidal category.  The PBW equivalence is the enveloping lane of chiral Koszul duality.  The mixed statement is the distributive-law theorem of Volume~I.
\end{proof}

This theorem replaces every dimension-based dispatch of an $\En$ level.  The geometry of the source selects the available directions; the operadic level records the operations actually constructed.

\chapter{Holomorphic Chern--Simons theory}
Let $X$ be a complex threefold with holomorphic volume form $\Omega_X$, and let $\mathfrak a$ be a quadratic Lie algebra. The BV fields are
\[
 \mathcal E_X=\Omega^{0,*}(X,\mathfrak a)[1],
\]
with degree $-1$ pairing
\[
 \omega(A,B)=\int_X\Omega_X\wedge\angles{A,B}.
\]
The action is
\[
 S_{\hCS}(A)=\int_X\Omega_X\wedge
 \left(\frac12\angles{A,\dbar A}+\frac16\angles{A,[A,A]}\right).
\]

\begin{theorem}[Classical hCS factorization algebra]
The action $S_{\hCS}$ satisfies the classical master equation. Its classical observables form a holomorphic factorization algebra on $X$.
\end{theorem}
\begin{proof}
The Hamiltonian vector field is the Maurer--Cartan differential $\dbar A+\frac12[A,A]$. Its square is the Dolbeault square, the derivation identity, and the Jacobi identity. The elliptic Dolbeault complex then gives factorization observables by the preceding theorem.
\end{proof}

A perturbative quantization is controlled by the local BV deformation complex. In complex dimension three, its one-loop gauge anomaly occupies the degree-four invariant-polynomial summand; mixed and gravitational components enter when the background fields carry the corresponding couplings. A quantization is a trivialization of the complete local anomaly class, compatible with restriction and factorization \cite{CostelloLiAnomaly}.

\chapter{Pushforward to a curve}
Let $p:X\to C$ be a holomorphic map to a complex curve. For a factorization algebra $\F$ on $X$, define
\[
 (p_*\F)(U)=\F(p^{-1}U).
\]

\begin{theorem}[Factorization pushforward]
The assignment $p_*\F$ is a factorization algebra on $C$. Suppose in addition that $p$ is proper and locally trivial in a topology compatible with the factorization theory, that $\F$ is locally constant in the fiber directions, and that product-chart excision and base change hold. Then on every trivializing open $U\subset C$ with fiber $F_p$ there is a natural equivalence
\[
 (p_*\F)(U)\simeq
 \int_{F_p}\F|_{F_p\times U},
\]
compatible on overlaps. If the resulting object is holomorphic and regular holonomic along $C$, Riemann--Hilbert algebraization gives a factorization $D$-module on $C$.
\end{theorem}
\begin{proof}
Disjoint opens in $C$ have disjoint inverse images, so the factorization products push forward. Weiss descent is inherited from inverse images of Weiss covers. On a product chart, local constancy in the fiber and factorization excision identify the value with fiberwise factorization homology. The base-change hypothesis glues these local equivalences. Regular-holonomic Riemann--Hilbert transports the holomorphic factorization object to the algebraic $D$-module setting.
\end{proof}

Thus a curve shadow is attached to an actual map and an actual field theory. Different fibrations and boundary conditions produce different ordered chiral algebras from the same categorical source.

\chapter{The derived centre and the physical bulk}
For an ordered boundary algebra $A$, define
\[
 Z_{\ord}(A)=\RHom_{A\tensor A^{\op}}(A,A).
\]
It is a canonical $\Etwo$ algebra. A bulk theory acting on the boundary gives
\[
 \beta:\Obs_{\mathrm{bulk}}|_{\partial}\longrightarrow Z_{\ord}(A).
\]

\begin{definition}[Open--closed completeness]
A boundary condition is open--closed complete on a protected sector when $\beta$ is an equivalence on that sector.
\end{definition}

The definition turns the phrase ``bulk equals centre'' into a comparison theorem. The centre exists algebraically for every ordered boundary algebra; the physical bulk is identified with it by the open--closed map.

\part{Hall and BPS algebras}

\chapter{Extension correspondences}
Let $\M=\coprod_{\gamma\in\Gamma}\M_\gamma$ be a moduli stack in an exact or derived category. Let $\M^{\mathrm{ext}}_{\alpha,\beta}$ classify extensions
\[
 0\longrightarrow E_\beta\longrightarrow E\longrightarrow E_\alpha\longrightarrow0.
\]
The correspondence is
\[
 \M_\alpha\times\M_\beta
 \xleftarrow{\ p\ }
 \M^{\mathrm{ext}}_{\alpha,\beta}
 \xrightarrow{\ q\ }
 \M_{\alpha+\beta}.
\]

\begin{theorem}[Hall associativity]
Assume a coefficient theory in which $p^*$ and $q_*$ are defined and satisfy base change for the two-step filtration stack. Then
\[
 a*b=q_*p^*(a\boxtimes b)
\]
defines an associative product.
\end{theorem}
\begin{proof}
The two parenthesizations are the two projections from the stack of filtrations $0\subset E_1\subset E_2\subset E_3$. Base change identifies both composites with pullback to this stack followed by pushforward to $E_3$.
\end{proof}

For critical CoHAs, the coefficient theory includes vanishing cycles and a coherent orientation. For symplectic CoHAs, it includes the selected equivariant Borel--Moore theory. Completion is chosen so that each output degree receives a convergent sum.

\chapter{Orientation data and the threefold Hall product}
For compact Calabi--Yau threefolds, natural orientation data on moduli of coherent sheaves and perfect complexes is constructed by Joyce and Upmeier \cite{JoyceUpmeier}.  A general cohomological Hall algebra for a $3$-Calabi--Yau category uses strong orientation data, which carries the coherent multiplicativity required by the attractor and extension correspondences \cite{KPS}.

\begin{theorem}[Threefold Hall algebra from strong orientation]
Let $\Ccal$ be a $3$-Calabi--Yau category satisfying the geometric hypotheses of Kinjo--Park--Safronov, and assume strong orientation data.  Then its cohomological Donaldson--Thomas complexes carry an associative cohomological Hall product.  The product is the parabolic-induction correspondence supplied by the attractor Lagrangian correspondence.
\end{theorem}
\begin{proof}
Kinjo--Park--Safronov prove the functoriality of the Donaldson--Thomas perverse sheaf under the attractor Lagrangian correspondence and use strong orientation to define the parabolic-induction morphism.  Their base-change and associativity theorem gives the CoHA \cite{KPS}.
\end{proof}

\begin{definition}[Reduced $K3\times E$ Hall datum]
A reduced $K3\times E$ Hall datum is a strong-oriented threefold Hall datum together with descent along the elliptic translation action, a reduced deformation theory, and a separated completion in physical charge.  A full quantum-group datum additionally includes a coproduct, Cartan extension, negative half, and nondegenerate continuous Hopf pairing.
\end{definition}

\chapter{The geometric \texorpdfstring{$E$}{E}-equivariant Hall source}
Let $S$ be a projective K3 surface, let $E$ be an elliptic curve, and put
\[
 X=S\times E.
\]
Write $\Coh_{\leq1}(X)$ for coherent sheaves whose support has dimension at most one.  This is an extension-closed abelian category.  Its numerical charge group is denoted $\Gamma_X$, and the degree of the projected one-cycle on $E$ is the additive function
\[
 m:\Gamma_X\longrightarrow\Z_{\geq0}.
\]
For $\gamma\in\Gamma_X$, let $\M_\gamma$ be the moduli stack of objects of charge $\gamma$.  Translation on the elliptic factor gives an action of $E$ on every $\M_\gamma$ and on every extension stack.

Oberdieck and Shen construct the $E$-equivariant Grothendieck ring of stacks and its Hall convolution.  In the present notation the completed algebra is
\[
 \Hcal_X^{E,\mathrm{mot}}
 =\widehat\bigoplus_{\gamma\in\Gamma_X}
 K_0^E(\operatorname{St}/\M_\gamma),
\]
where the completion is taken in effective curve class and Euler characteristic.

\begin{theorem}[Geometric reduced Hall algebra]
The extension correspondence in $\Coh_{\leq1}(X)$ defines an associative $E$-equivariant motivic Hall product on $\Hcal_X^{E,\mathrm{mot}}$.  The Oberdieck--Shen reduced integration morphism
\[
 I_X^E:\Hcal_X^{E,\mathrm{mot}}
 \longrightarrow
 \widehat{\Q[\Gamma_X]}[\epsilon]/(\epsilon^2)
\]
is multiplicative on the regular subalgebra on which their integration theorem is defined.  The coefficient of $\epsilon$ is the reduced contribution of the positive-dimensional $E$-orbits.
\end{theorem}
\begin{proof}
The stack of short exact sequences is $E$-equivariant because translation preserves exactness and numerical charge.  The stack of two-step filtrations identifies the two composites of three Hall products, so equivariant base change proves associativity.  The integration statement is the $E$-equivariant motivic Hall theorem of Oberdieck--Shen \cite{OberdieckShen}.  Their stabilizer stratification separates fixed contributions from moving elliptic orbits and packages the latter in the square-zero parameter $\epsilon$.
\end{proof}

\chapter{Local and wrapped geometry inside the Hall algebra}
Define the completed degree-zero and positive-degree pieces
\[
 \Hcal_X^{\mathrm{loc}}
 =\widehat\bigoplus_{m(\gamma)=0}
 K_0^E(\operatorname{St}/\M_\gamma),
 \qquad
 \Hcal_X^{\mathrm{wr}}
 =\widehat\bigoplus_{m(\gamma)>0}
 K_0^E(\operatorname{St}/\M_\gamma).
\]
A pure one-dimensional support of degree zero projects to a finite subset of $E$.  A support of positive degree projects onto $E$.  The first sector is point-local on the elliptic curve; the second is wrapped.

\begin{theorem}[Geometric local--wrapped decomposition]
Hall multiplication respects elliptic degree:
\[
 \Hcal_{m_1}*\Hcal_{m_2}\subseteq\Hcal_{m_1+m_2}.
\]
Consequently $\Hcal_X^{\mathrm{loc}}$ is a subalgebra, $\Hcal_X^{\mathrm{wr}}$ is a closed two-sided ideal, and
\[
 \Hcal_X^{E,\mathrm{mot}}
 \cong
 \Hcal_X^{\mathrm{loc}}\ltimes\Hcal_X^{\mathrm{wr}}
\]
as a complete vector space with multiplication determined by the Hall correspondences.  The reduced integration morphism preserves this decomposition.
\end{theorem}
\begin{proof}
The cycle of the middle term of a short exact sequence is the sum of the cycles of its subobject and quotient.  Projection to $E$ therefore adds degrees.  Degree zero is closed under addition, and positive degree absorbs addition by a nonnegative degree.  The same grading is present on the target group algebra of $I_X^E$.
\end{proof}

This theorem supplies the geometric version of the local--wrapped splitting constructed from the Borcherds root algebra in the companion volume.  The two constructions have the same additive elliptic degree and become comparable after a primitive root framing is chosen.

\chapter{Cohomological refinement and the Gram character}
A cohomological refinement uses the Donaldson--Thomas sheaf on the derived moduli stack.  The datum required for the convolution is the strong orientation coherence of Kinjo--Park--Safronov.  A strong orientation enhancement of the natural Joyce--Upmeier orientation therefore produces
\[
 \Hcal_X^{E,\mathrm{coh}}
 =\widehat\bigoplus_{\gamma\in\Gamma_X}
 H_*^{\mathrm{BM},E}(\M_\gamma,\Phi_\gamma),
\]
with its cohomological Hall product.

\begin{theorem}[Strong-oriented cohomological lift]
For every strong orientation enhancement on the $E$-equivariant derived moduli problem, the attractor correspondence constructs an associative cohomological Hall algebra $\Hcal_X^{E,\mathrm{coh}}$.  Elliptic degree again gives a subalgebra in degree zero and a closed ideal in positive degree.  Forgetting the cohomological coefficient system recovers the corresponding motivic Hall correspondences.
\end{theorem}
\begin{proof}
The parabolic-induction theorem of Kinjo--Park--Safronov gives the cohomological product from the strong-oriented attractor Lagrangian correspondence \cite{KPS}.  Numerical charge and elliptic degree are constant on connected components of every correspondence, so the proof of the preceding decomposition applies verbatim.
\end{proof}

For the dyonic charge lattice
\[
 \Gamma_{\mathrm{phys}}
 =\widetilde H(S,\Z)\oplus\widetilde H(S,\Z),
\]
write
\[
 \Pi(Q,P)=\left(\frac{Q^2}{2},Q\cdot P,\frac{P^2}{2}\right).
\]
To compare sheaf charges with this lattice, choose an additive map
\[
 c:\Gamma_X\longrightarrow\Gamma_{\mathrm{phys}}.
\]
Put $B_\Pi(a,b)=\Pi(a+b)-\Pi(a)-\Pi(b)$ and give
$\widetilde\Gamma_{\mathrm{phys}}=\Gamma_{\mathrm{phys}}\times\mathbb Z^3$
the group law
\[
 (a,u)(b,v)=(a+b,u+v+B_\Pi(a,b)).
\]
Then $j_\Pi(a)=(a,\Pi(a))$ is a homomorphism.
Assume $j_\Pi c$ preserves the selected charge completions.
On the regular subalgebra $\mathcal H_{\mathrm{reg}}$ of the reduced
integration theorem, define
\[
 I_{X,\Pi,c}^E=(j_\Pi c)_*I_X^E:
 \mathcal H_{\mathrm{reg}}\longrightarrow
 \widehat{\mathbb Q[\widetilde\Gamma_{\mathrm{phys}}]}
 [\epsilon]/(\epsilon^2).
\]
Write $I_{X,\Pi}^E$ for this map when $c$ has been fixed.

\begin{proposition}[Polarized reduced character]\label{cyiv:polarized-character}
The map $I_{X,\Pi,c}^E$ is multiplicative.
Its coefficients are the reduced invariants supplied by $I_X^E$,
regraded by $j_\Pi c$.
If $g\in1+F^1H$ is group-like in a complete filtered Hall bialgebra
and the logarithm converges, then $\log g$ is primitive.
\end{proposition}
\begin{proof}
Expanding $\Pi(a+b+d)$ gives
\[
 B_\Pi(a,b)+B_\Pi(a+b,d)=B_\Pi(b,d)+B_\Pi(a,b+d).
\]
This proves associativity of the group law.
The definition gives $j_\Pi(a+b)=j_\Pi(a)j_\Pi(b)$.
Composing its group-algebra map with reduced integration proves
multiplicativity on $\mathcal H_{\mathrm{reg}}$.
For the last assertion, the two factors in
$\Delta g=(g\widehat\otimes1)(1\widehat\otimes g)$ commute.
Taking the convergent logarithm gives
$\Delta\log g=\log g\widehat\otimes1+1\widehat\otimes\log g$.
\end{proof}

The primitive reduced Gromov--Witten theorem supplies the separate
character $-\chi_{10}^{-1}$ in classes primitive in the K3 factor
\cite{OP}.
Comparing this character with the Hall character requires the same
charge map, sector, and normalization.
A primitive root framing must also preserve the extension bracket
and parity of the chosen BPS states.

\chapter{BPS primitives}
A connected graded bialgebra has primitive Lie algebra
\[
 \Prim(H)=\set{x:\Delta x=x\tensor1+1\tensor x}.
\]
In cohomological Donaldson--Thomas theory, the BPS sheaf isolates the primitive part of the Hall algebra.

\begin{theorem}[Two-Calabi--Yau BPS algebra]
For the geometric two-Calabi--Yau categories covered by the BPS-algebra theorem, the BPS algebra is the positive enveloping algebra of a generalized Kac--Moody Lie algebra. This scope includes semistable sheaves on K3 and other symplectic surfaces under the stated good-moduli hypotheses.
\end{theorem}
\begin{proof}
The theorem is the Davison--Hennecart--Schlegel Mejia construction \cite{DHS}. The simple BPS sheaves give the simple-root spaces, extension correspondences give the bracket, and the integrality/PBW theorem identifies the Hall algebra with the positive enveloping algebra.
\end{proof}

This result supplies a global K3 generalized Kac--Moody algebra at the two-Calabi--Yau Hall level. Yangian structures arise from additional equivariant correspondences and coproducts in the surface and quiver models described below.

\chapter{The integral Hall--Borcherds amalgam}
Let $\mathfrak p_X^+$ be a complete charge-graded BPS Lie superalgebra supplied by a geometric Hall source, and let $\nDelta$ be the positive Igusa Borcherds Lie superalgebra.  The polarized charge extension places both gradings over one abelian Cartan algebra $\mathfrak h_\Pi$.  Form the two complete Borels
\[
 \mathfrak b_X=\mathfrak h_\Pi\ltimes\mathfrak p_X^+,
 \qquad
 \mathfrak b_\Delta=\mathfrak h_\Pi\ltimes\nDelta.
\]
The Cartan acts through the physical charge on the first Borel and through Gram degree on the second.

\begin{definition}[Integral protected Lie algebra]
The integral Hall--Borcherds algebra is the completed amalgamated coproduct
\[
 \mathfrak b_{X\boxtimes\Delta}
 =\mathfrak b_X
  \widehat\amalg_{\mathfrak h_\Pi}
  \mathfrak b_\Delta
\]
in complete charge-filtered dg Lie superalgebras.
\end{definition}

\begin{theorem}[Hall--Borcherds amalgamation]
The completed amalgam exists and represents pairs of continuous Lie morphisms from $\mathfrak b_X$ and $\mathfrak b_\Delta$ that agree on $\mathfrak h_\Pi$.  The Cartan projections of the two Borels induce canonical retractions
\[
 r_X:\mathfrak b_{X\boxtimes\Delta}\longrightarrow\mathfrak b_X,
 \qquad
 r_\Delta:\mathfrak b_{X\boxtimes\Delta}\longrightarrow\mathfrak b_\Delta.
\]
Its completed enveloping algebra is the amalgamated free product
\[
 \widehat U(\mathfrak b_{X\boxtimes\Delta})
 \cong
 \widehat U(\mathfrak b_X)
 \widehat *_{\widehat U(\mathfrak h_\Pi)}
 \widehat U(\mathfrak b_\Delta).
\]
The geometric Hall operations survive under $r_X$, while the Igusa bracket, PBW character, and automorphic determinant survive under $r_\Delta$.
\end{theorem}
\begin{proof}
Construct the ordinary Lie coproduct from the free product of the two Borels and divide by the ideal identifying their two Cartan copies.  Complete by total positive charge and proper root height.  Every fixed filtered degree uses finitely many Lie words, so the completion is separated and has the stated mapping property.  The projection $\mathfrak b_X\to\mathfrak h_\Pi\to\mathfrak b_\Delta$ together with the identity on $\mathfrak b_\Delta$ gives $r_\Delta$ by universality; the second retraction is symmetric.  The enveloping functor preserves coproducts of Lie algebras, and PBW identifies the completed algebraic filtration with the displayed completed free product.
\end{proof}

\begin{definition}[Comparison ideal]
A primitive comparison $\varphi:\nDelta\to\mathfrak p_X^+$ preserving Cartan degree defines the closed ideal
\[
 \mathcal J_\varphi
 =\overline{\angles{
   i_\Delta(x)-i_X(\varphi(x))\mid x\in\nDelta
 }}
 \subset\mathfrak b_{X\boxtimes\Delta}.
\]
The diagonal realization is
\[
 \mathfrak b_{X\simeq\Delta,\varphi}
 =\mathfrak b_{X\boxtimes\Delta}/\mathcal J_\varphi.
\]
\end{definition}

\begin{proposition}[Graph quotient of a primitive comparison]\label{cyiv:graph-quotient}
Let $\varphi:\mathfrak b_\Delta\to\mathfrak b_X$ be a continuous Lie
morphism restricting to the identity on $\mathfrak h_\Pi$.
Use a completion in which the amalgam and closed quotient have
their stated universal properties.
Then
\[
 \mathfrak b_{X\boxtimes\Delta}/\mathcal J_\varphi
 \cong\mathfrak b_X.
\]
The two factor maps become $\id$ and $\varphi$.
If $\varphi$ is an isomorphism, the quotient is also isomorphic
to $\mathfrak b_\Delta$.
\end{proposition}
\begin{proof}
The maps $\id$ and $\varphi$ agree on $\mathfrak h_\Pi$.
They give $q_\varphi:\mathfrak b_{X\boxtimes\Delta}\to\mathfrak b_X$,
which kills the generators of $\mathcal J_\varphi$.
The inclusion of $\mathfrak b_X$ gives a map $s$ in the other direction.
The composite $q_\varphi s$ is the identity.
The other composite is the identity on the geometric factor.
On the Borcherds factor, the defining relation identifies $x$ with
$s\varphi(x)$, so it is also the identity there.
The two factor images topologically generate the quotient.
Continuity proves that the maps are inverse on the completion.
\end{proof}

The graph quotient implements a specified primitive comparison.
Constructing that comparison requires geometric simple states whose
extension brackets satisfy the Borcherds relations.
The quotient formula alone produces neither those states nor their
orientation data.

\chapter{Coproducts, pairings, and doubles}
A positive Hall algebra becomes a quantum group after additional structure is supplied. Let $H^+$ be a complete graded bialgebra, let $H^0$ be a Cartan algebra, and let $H^-$ be the opposite or dual half. Let
\[
 \angles{-,-}:H^+\widehattensor H^-\longrightarrow\kfield
\]
be a nondegenerate Hopf pairing.

\begin{definition}[Completed Hall double]
The completed Drinfeld double is the topological algebra
\[
 \widehat D(H^+)=H^-\widehattensor H^0\widehattensor H^+
\]
with internal products in the three factors and cross-relations determined by the Hopf pairing. Its representation category is monoidal by the coproduct and braided by the universal $R$-matrix whenever the canonical tensor converges.
\end{definition}

\begin{theorem}[Double construction]
Let $H^+$ and $H^-$ be paired complete graded bialgebras, let $H^0$ be a Cartan bialgebra acting by the stated weight grading, and let
\[
 \angles{-,-}:H^+\widehattensor H^-\longrightarrow\kfield
\]
be a continuous nondegenerate Hopf pairing. Assume the triangular multiplication map from $H^-\widehattensor H^0\widehattensor H^+$ is separated and complete. Then the pairing cross-relations determine a unique completed double with this triangular decomposition. On every continuous module category on which the canonical tensors converge, the completed canonical element gives the universal $R$-matrix.
\end{theorem}
\begin{proof}
The Hopf pairing defines the cross-relations between the paired halves and makes the three canonical embeddings bialgebra morphisms. The Hopf identities identify the two reductions of every triple product, proving associativity. Separated triangular decomposition gives uniqueness. Finite-degree canonical tensors form a compatible inverse system; convergence on the selected module category yields the quasitriangular element \cite{DrinfeldDouble,EtingofSchiffmann}.
\end{proof}

\begin{theorem}[Quantum Hall--Borcherds amalgam]
Let $H_X$ and $H_\Delta$ be complete Hopf algebras containing the same complete Cartan Hopf algebra $H^0$, with continuous Hopf retractions onto $H^0$.  Their completed amalgamated free product
\[
 H_{X\boxtimes\Delta}
 =H_X\widehat *_{H^0}H_\Delta
\]
has a unique complete Hopf structure extending the two coproducts, counits, and antipodes.  It represents pairs of continuous Hopf morphisms agreeing on $H^0$ and has canonical Hopf retractions to both factors.  A mixed universal $R$-matrix is equivalent to a continuous cross-pairing satisfying the two hexagon identities; the intrinsic $R$-matrices of the two retracts remain defined independently.
\end{theorem}
\begin{proof}
The algebraic amalgamated free product exists because complete filtered algebras are obtained degreewise from finite quotients.  The two coproducts agree on $H^0$, so the universal property extends them to a coproduct on the free product.  Coassociativity is checked on the two generating Hopf subalgebras.  The same argument extends the counit and antipode.  The Hopf retractions are induced by the factor retractions onto $H^0$.  A quasitriangular structure extending both factors requires the additional commutation data between the two generating halves; writing this data as a cross-pairing turns the quasitriangular and Yang--Baxter equations into the two hexagon identities.
\end{proof}

The deconcatenation coproduct on a bar complex and the representation-theoretic coproduct on a quantum group occupy different structures. A geometric comparison between them is a theorem about a selected model.

\chapter{Ordered chiral bialgebras}
A tensor product of representations requires compatible actions on
threefold products and the unit.
Retain the factorization coefficient category
$(\Fact_D(C),\tensor^!)$ on the smooth complex curve $C$.
Fix a category $\widehat{\mathcal V}_C$ of complete filtered objects
with specified presentations $F\simeq\varprojlim_r F_r$,
where $F_r\in\Fact_D(C)$.
Use derived limits in the fixed Ran enhancement and define
\[
 F\widehat{\tensor}^{!}G
 =\varprojlim_{r,s}(F_r\tensor^!G_s).
\]
Assume these limits remain factorizing and give
$\widehat{\mathcal V}_C$ a symmetric monoidal structure with unit $\mathbf1$.
Its associativity, symmetry, and unit equivalences must preserve
the presentations and the Ran structure.
Morphisms are continuous maps compatible with these presentations.
These are hypotheses on the chosen completion.

\begin{definition}[Ordered chiral bialgebra]\label{def:vol4-ordered-chiral-bialgebra}
An ordered chiral bialgebra is a counital $\Eone$-coalgebra object
in complete ordered chiral algebras:
\[
 A\in\operatorname{CoAlg}_{\Eone}
 \left(\Alg_{\Eone}
 (\widehat{\mathcal V}_C,\widehat{\tensor}^{!})\right).
\]
It has continuous multiplication $m$ and unit $\eta$, and continuous
morphisms of unital ordered chiral algebras
\[
 \Delta:A\longrightarrow A\widehat{\tensor}^{!}A,
 \qquad \epsilon:A\longrightarrow\mathbf1.
\]
The counit is the chosen augmentation.
Write $\tau$ for the coefficient symmetry, including Koszul signs.
Suppressing ambient associativity and unit equivalences, compatibility reads
\[
\begin{aligned}
 \Delta m&=(m\widehat{\tensor}^{!}m)
  (\id\widehat{\tensor}^{!}\tau\widehat{\tensor}^{!}\id)
  (\Delta\widehat{\tensor}^{!}\Delta),\\
 \Delta\eta&=\eta\widehat{\tensor}^{!}\eta,
 \qquad \epsilon m=\epsilon\widehat{\tensor}^{!}\epsilon,
 \qquad \epsilon\eta=\id_{\mathbf1}.
\end{aligned}
\]
Coassociativity and the counit identities read
\[
\begin{aligned}
 (\Delta\widehat{\tensor}^{!}\id)\Delta
   &=(\id\widehat{\tensor}^{!}\Delta)\Delta,\\
 (\epsilon\widehat{\tensor}^{!}\id)\Delta
   &=\id_A=(\id\widehat{\tensor}^{!}\epsilon)\Delta.
\end{aligned}
\]
For a homotopy-coherent presentation, these identities include the
higher coherences of the displayed coalgebra object.
Every operation and coherence preserves diagonal transitions and
disjoint-locus factorization maps.
An antipode is a continuous convolution inverse $S:A\to A$ of $\id_A$:
\[
 m(S\widehat{\tensor}^{!}\id)\Delta
 =\eta\epsilon
 =m(\id\widehat{\tensor}^{!}S)\Delta.
\]
A bialgebra with an antipode is a Hopf object in this coefficient category.
\end{definition}

A continuous left $A$-module is an object $M\in\widehat{\mathcal V}_C$
with a unital $\Eone$ action
$\rho_M:A\widehat{\tensor}^{!}M\to M$.
Its action and coherences are morphisms in these factorization coefficients.
Write $\operatorname{Mod}_A(\widehat{\mathcal V}_C)$ for their category.

\begin{theorem}[Representation fusion]\label{thm:vol4-representation-fusion}
An ordered chiral bialgebra makes
$\operatorname{Mod}_A(\widehat{\mathcal V}_C)$ monoidal.
Its product is $M\widehat{\tensor}^{!}N$ with action
\[
\begin{aligned}
 A\widehat{\tensor}^{!}(M\widehat{\tensor}^{!}N)
 &\xrightarrow{\Delta\widehat{\tensor}^{!}\id}
 (A\widehat{\tensor}^{!}A)\widehat{\tensor}^{!}
 (M\widehat{\tensor}^{!}N)\\
 &\xrightarrow{\id\widehat{\tensor}^{!}\tau_{A,M}
                       \widehat{\tensor}^{!}\id}
 (A\widehat{\tensor}^{!}M)\widehat{\tensor}^{!}
 (A\widehat{\tensor}^{!}N)\\
 &\xrightarrow{\rho_M\widehat{\tensor}^{!}\rho_N}
 M\widehat{\tensor}^{!}N.
\end{aligned}
\]
The unit is $\mathbf1$ with action induced by $\epsilon$.
The ambient associator and unit equivalences are module equivalences.

Suppose a continuous universal $R$-matrix is also given.
It is an invertible morphism
$R:\mathbf1\to A\widehat{\tensor}^{!}A$ for the tensor-product algebra,
satisfying
\[
\begin{gathered}
 R\Delta(a)=\Delta^{\op}(a)R,\qquad
 \Delta^{\op}=\tau_{A,A}\Delta,\\
 (\Delta\widehat{\tensor}^{!}\id)(R)=R_{13}R_{23},
 \qquad
 (\id\widehat{\tensor}^{!}\Delta)(R)=R_{13}R_{12},\\
 (\epsilon\widehat{\tensor}^{!}\id)(R)=\eta
 =(\id\widehat{\tensor}^{!}\epsilon)(R).
\end{gathered}
\]
The first equation is an identity of morphisms from $A$ to its completed tensor square.
Leg notation uses the unit and symmetry of $\widehat{\mathcal V}_C$.
In a homotopy-coherent presentation, require compatible lifts of these identities.
Then $c_{M,N}=\tau_{M,N}\circ R_{M,N}$ is a braiding,
where $R_{M,N}$ acts through the external tensor-product module action.
\end{theorem}
\begin{proof}
The symmetry combines the independent module actions into an
$A\widehat{\tensor}^{!}A$-action.
Restriction along the unital algebra map $\Delta$ gives the displayed action.
For three modules, the two actions use
$(\Delta\widehat{\tensor}^{!}\id)\Delta$ and
$(\id\widehat{\tensor}^{!}\Delta)\Delta$.
Coassociativity makes the ambient associator $A$-linear.
The counit identities make the ambient unit equivalences $A$-linear.
The coalgebra coherences give the pentagon and triangle coherences.
The completion hypotheses keep these maps continuous and factorizing.
Quasitriangularity makes $c_{M,N}$ an intertwiner and invertibility makes it an equivalence.
The two coproduct identities for $R$ give the hexagons.
Its counit identities give compatibility with the unit.
\end{proof}

The tensor in this construction is the completed transverse tensor
$\tensor^!$ on factorization coefficients.
A chiral-convolution or Hall coproduct requires its own construction
and a comparison in the specified category.
A spectral version additionally requires a difference coordinate,
a Laurent-series completion, and compatible domains for iterated expansions.
Its identities must hold in those domains before defining module actions.

\begin{example}[A counital algebra map without coassociativity]
In ordinary complex vector spaces, let $A=\C[x]$ and put
\[
 \Delta(x)=x\tensor1+1\tensor x+x\tensor x^2,
 \qquad \epsilon(x)=0.
\]
Polynomial substitution extends these assignments to unital algebra maps.
Writing $F(u,v)=u+v+uv^2$, the equations
$F(0,v)=v$ and $F(u,0)=u$ prove both counit identities.
But
\[
\begin{split}
 F(F(u,v),w)-F(u,F(v,w))
 &=-uv^2w^4-uv^2w^2-2uvw^3-2uvw\\
 &=-uvw(1+w^2)(vw+2)\ne0.
\end{split}
\]
On three regular modules, the two actions differ already on
$1\tensor1\tensor1$.
Thus the canonical tensor associator is not $A$-linear.
This ordinary example shows why multiplicativity and the counit identities
must be supplemented by coassociativity.
It makes no assertion about a geometric realization in factorization coefficients.
\end{example}

\part{Exact model cases}

\chapter{The \texorpdfstring{$A_2$}{A2} Hall algebra}
For the quiver $1\to2$ over $\mathbb F_q$, the twisted Hall algebra has generators $E_1,E_2$ and $v^2=q$. Counting subrepresentations gives
\[
 E_1^2E_2-(v+v^{-1})E_1E_2E_1+E_2E_1^2=0.
\]
For a representation of dimension vector $(2,1)$, the arrow has rank zero or one. The relevant flag counts are
\[
\begin{array}{c|ccc}
\rank&112&121&211\\ \hline
0&q+1&q+1&q+1\\
1&q+1&1&0.
\end{array}
\]
After the Hall twists $v^{-1},1,v$, both rows give zero. This is the elementary geometric origin of the quantum Serre relation \cite{Ringel,Green}.

\chapter{The tripled Jordan quiver and \texorpdfstring{$\mathbb C^3$}{C3}}
Let $Q$ have one vertex and loops $x,y,z$, with potential
\[
 W=\Tr x[y,z].
\]
The Jacobi algebra is $\C[x,y,z]$. The critical CoHA uses the vanishing cycles of $\Tr W$ on the representation spaces. Its positive currents are generated by the one-vertex dimension sectors.

\begin{theorem}[Affine-Yangian positive half]
After the standard equivariant localization and completion, the critical CoHA of the tripled Jordan quiver is identified with the positive half of the affine Yangian of $\mathfrak{gl}_1$. The independently defined affine Yangian carries its standard Cartan and negative current halves; the positive-half comparison embeds the CoHA into that triangular presentation.
\end{theorem}
\begin{proof}
The CoHA product becomes the shuffle product under localization. Dimensional reduction identifies the tripled-quiver critical theory with the equivariant surface CoHA, and Davison's affine-BPS theorem identifies its completed positive algebra with the strictly positive affine Yangian. The shuffle kernel satisfies the affine-Yangian current relations with equivariant parameters $\epsilon_1,\epsilon_2,\epsilon_3$ and Calabi--Yau relation $\epsilon_1+\epsilon_2+\epsilon_3=0$. PBW and current comparison identify the positive algebra. The independently presented full affine Yangian is the triangular algebra generated by its negative currents, Cartan currents, and this positive half \cite{DavisonAffineBPS,DavisonMeinhardt,SVYangians,Tsymbaliuk}.
\end{proof}

The vacuum representation produces a $W_{1+\infty}$-type vertex algebra through an evaluation homomorphism. The CoHA, its double, and the vacuum algebra occupy successive levels: positive algebra, paired quantum group, and distinguished representation.

\chapter{Surface CoHAs and ADE Yangians}
Let $S_\Gamma\to\C^2/\Gamma$ be the minimal resolution for a finite subgroup $\Gamma\subset\operatorname{SL}_2(\C)$. The exceptional curves have ADE intersection matrix. Moduli of sheaves supported on these curves carry convolution correspondences.

\begin{theorem}[Curve-supported surface CoHA]
The completed positive CoHA of one-dimensional sheaves supported on the exceptional curves is identified with a nonstandard positive half of the affine Yangian of the corresponding ADE Lie algebra.
\end{theorem}
\begin{proof}
The theorem is the curve-supported surface construction of Diaconescu--Porta--Sala--Schiffmann--Vasserot \cite{DPPSSV}. The simple curve sheaves give the simple generators, Hecke correspondences give the current operators, and the completed PBW comparison gives the affine Yangian presentation.
\end{proof}

This theorem is a local K3 model whenever an ADE configuration occurs on a K3 surface. The global K3 surface adds compact curve classes and monodromy beyond the local root lattice.

\chapter{Coulomb branches and shifted Yangians}
For a three-dimensional $\mathcal N=4$ gauge theory, the BFN Coulomb branch is defined by equivariant Borel--Moore homology of an affine-Grassmannian correspondence. Its convolution algebra quantizes the Coulomb branch \cite{BFN}. In quiver gauge theories, shifted Yangians and their truncations provide presentations of these quantizations. Boundary conditions select the shift and truncation data.

The BFN construction supplies a second geometric source of Yangians. Its grading and coproduct arise from loop rotation and convolution, while the surface CoHA source arises from sheaf extensions. A comparison between the two is a model-specific equivalence.

\part{K3 symmetries}

\chapter{The Mukai lattice}
For a K3 surface $S$, the Mukai lattice is
\[
 \Mukai(S,\Z)=H^0(S,\Z)\oplus H^2(S,\Z)\oplus H^4(S,\Z)
\]
with pairing
\[
 ((r,c,s),(r',c',s'))=c\cdot c'-rs'-r's.
\]
It is even unimodular of signature $(4,20)$ and rank $24$.

\chapter{The Nakajima Heisenberg algebra}
Let
\[
 \mathbb H_S=\bigoplus_{n\ge0}H^*(S^{[n]}).
\]
For $m\in\Z\setminus\{0\}$ and $\alpha\in H^*(S)$, the Nakajima correspondence gives an operator $p_m(\alpha)$.

\begin{theorem}[K3 Heisenberg action]
The operators satisfy
\[
 [p_m(\alpha),p_n(\beta)]
 =m\,\delta_{m+n,0}\int_S\alpha\cup\beta\;\id.
\]
They generate the rank-twenty-four Heisenberg algebra acting irreducibly on $\mathbb H_S$ with vacuum $1\in H^*(S^{[0]})$.
\end{theorem}
\begin{proof}
The commutator is the excess intersection of the two incidence correspondences. Opposite creation and annihilation correspondences meet along the diagonal with multiplicity $m$ and intersection pairing $\int_S\alpha\beta$. Correspondences of the same sign commute. The vacuum construction gives the Fock representation \cite{Nakajima97}.
\end{proof}

This action is global, compact, and abelian at the Lie level. It is an exact K3 symmetry and supplies the rank-twenty-four Fock sector of any broader K3 algebra.

\chapter{The orthogonal LLV algebra}
For a compact hyperkähler manifold, every Lefschetz class gives an $\mathfrak{sl}_2$ triple on cohomology. The Lie algebra generated by all such triples is the Looijenga--Lunts--Verbitsky algebra.

\begin{theorem}[Orthogonal LLV symmetries]
For a K3 surface $S$, the LLV algebra acting on $H^*(S)$ is
\[
 \mathfrak g_{\mathrm{LLV}}(S)
 \cong\mathfrak{so}\bigl(\widetilde H(S),(-,-)_{\mathrm{Muk}}\bigr)
 \cong\mathfrak{so}(4,20).
\]
For an irreducible hyperk\"ahler manifold $M$ of $K3^{[n]}$ type with $n\ge2$, one has $b_2(M)=23$ and
\[
 \mathfrak g_{\mathrm{LLV}}(M)\cong\mathfrak{so}(4,21).
\]
Moduli spaces of stable objects on $S$ carry their own LLV actions; Mukai and monodromy correspondences relate these actions to the K3 lattice.
\end{theorem}
\begin{proof}
The Looijenga--Lunts--Verbitsky theorem identifies the Lie algebra generated by all Lefschetz triples with the orthogonal algebra of the Mukai extension $H^2(M)\oplus U$. For $S$, the extension has signature $(4,20)$. For $K3^{[n]}$ type with $n\ge2$, the Beauville--Bogomolov lattice has signature $(3,20)$, so adjoining the hyperbolic plane gives signature $(4,21)$ \cite{LooijengaLunts,Verbitsky96}.
\end{proof}

The LLV algebra is a classical cohomological symmetry. A Yangian deformation of the full orthogonal action consists of two further structures: a Drinfeld current presentation and a compatible spectral coproduct.

\chapter{Local ADE Yangians}
An ADE configuration of $(-2)$-curves on $S$ gives a root sublattice of $H^2(S,\Z)$. Sheaves supported on the configuration produce the surface CoHA of the preceding part and hence a completed positive affine ADE Yangian. Nakajima quiver varieties supply compatible Yangian actions on their equivariant cohomology.

Thus K3 geometry contains local nonabelian Yangian sectors indexed by ADE enhancements of the Picard lattice. Picard monodromy transports the ADE root sublattices over the corresponding Noether--Lefschetz loci. A family of sheaf moduli and Hecke correspondences compatible with that monodromy transports the Yangian actions themselves.

\chapter{The K3 symmetry trichotomy}
The three exact structures are
\begin{center}
\small
\begin{tabularx}{0.96\textwidth}{@{}>{\raggedright\arraybackslash}p{0.23\textwidth}>{\raggedright\arraybackslash}X>{\raggedright\arraybackslash}p{0.28\textwidth}@{}}
\toprule
\textbf{Structure}&\textbf{Carrier}&\textbf{Operation}\\
\midrule
Mukai Heisenberg&$\bigoplus_nH^*(S^{[n]})$&creation and annihilation\\
LLV $\mathfrak{so}(4,20)$&$H^*(S)$; separately $\mathfrak{so}(4,21)$ on $K3^{[n]}$ type&Lefschetz Lie action\\
Local ADE Yangian&curve-supported equivariant moduli&Hecke and current correspondences\\
\bottomrule
\end{tabularx}
\end{center}
A global K3 quantum algebra is a system that contains these three actions with explicit comparison homomorphisms. The two-Calabi--Yau BPS theorem supplies a global generalized Kac--Moody positive algebra. A full Yangian or double is specified by a compatible coproduct, spectral parameter, negative half, Cartan part, pairing, and completion.

\chapter{K3 correspondences on fixed coefficients}
Fix rational cohomology and a collection $\mathscr M_S$ of smooth
projective moduli spaces on a projective K3 surface $S$.
Include the Hilbert schemes $S^{[n]}$ and close the collection under
finite products.
For $M,N\in\mathscr M_S$, define the graded kernel space
\[
 K(M,N)=H^*(M\times N,\mathbb Q)[2\dim_{\mathbb C}M].
\]
A kernel of degree $r$ acts from $H^j(M)$ to $H^{j+r}(N)$ by
\[
 T_\alpha(v)=(p_N)_*(p_M^*v\cup\alpha).
\]
The shift in $K(M,N)$ accounts for integration over $M$.
All pushforwards in this construction are proper.

\begin{proposition}[Kernel composition]\label{cyiv:k3-kernels}
The product
\[
 \beta\circ\alpha=(p_{13})_*
 (p_{12}^*\alpha\cup p_{23}^*\beta)
\]
defines an associative graded category with diagonal identities.
The kernel action identifies $K(M,N)$ with
$\operatorname{Hom}^*(H^*(M),H^*(N))$.
\end{proposition}
\begin{proof}
The K\"unneth isomorphism and Poincar\'e duality identify
$H^*(M\times N)[2\dim M]$ with $H^*(M)^\vee\otimes H^*(N)$.
Under this identification, $T_\alpha$ is the associated linear map.
Integration over the middle factor contracts its dual pair, so
$T_{\beta\circ\alpha}=T_\beta T_\alpha$.
Faithfulness then proves associativity and the diagonal identity.
External products give tensor products with their Koszul signs.
\end{proof}

A geometric symmetry is specified by selected kernels in this category.
Nakajima incidence kernels connect different Hilbert schemes.
Lefschetz multiplication and its adjoint act on each individual
cohomology space.
Using every cohomology kernel would instead give all linear maps;
the selected geometric kernels carry the representation-theoretic content.

\chapter{The algebra generated by geometric actions}
Put $V_S=\bigoplus_{M\in\mathscr M_S}H^*(M,\mathbb Q)$.
Choose a set of homogeneous operators $T_\lambda$ on $V_S$, each sending
a vector with finite support to a vector with finite support.
The Nakajima operators on the Hilbert-scheme summand satisfy this
condition: each changes the point number by a fixed integer.
Individual Lefschetz operators, extended by zero on other summands,
also satisfy it.
Let $F$ be the free unital graded algebra on symbols $t_\lambda$.

\begin{proposition}[Geometric action algebra]\label{cyiv:k3-action}
The map
\[
 \rho:F\longrightarrow\operatorname{End}_{\mathbb Q}(V_S),
 \qquad t_\lambda\longmapsto T_\lambda
\]
defines a faithful action of $A_S=F/\ker\rho$.
For every algebra $B$, maps $A_S\to B$ are exactly maps $F\to B$
that kill $\ker\rho$.
Products and commutators of the selected kernels are their convolution
products and commutators on the relevant moduli spaces.
\end{proposition}
\begin{proof}
Finite words in the selected operators preserve finite support, so
$\rho$ is defined by composition.
The quotient acts faithfully by definition of its kernel.
The factorization property is the universal property of an algebra
quotient.
The convolution assertion follows from Proposition~\ref{cyiv:k3-kernels}.
\end{proof}

An equivariant ADE construction uses a separate coefficient field,
for example $\operatorname{Frac}H_T^*(\mathrm{pt},\mathbb Q)$,
and its prescribed localization and proper pushforwards.
Joining it to $A_S$ requires maps between the actual state spaces that
intertwine those operators and their coefficients.
A completion of $A_S$ requires a specified topology and continuity of
every generator.
The indefinite Mukai norm alone supplies neither condition.
The desired common K3 quantum algebra must construct these local-to-global
actions and their mixed convolution relations.

\chapter{Tensor reconstruction and geometric kernels}
Let $\mathscr C$ be an essentially small rigid linear monoidal category
and let $\omega:\mathscr C\to\operatorname{Vect}_{\mathrm{fd}}(\mathbb Q)$
be a strong monoidal functor.
A choice of geometric kernels specifies the morphisms of $\mathscr C$;
it is part of the input.
Define
\[
 H=\left(\bigoplus_{M\in\mathscr C}\omega(M)^*\otimes\omega(M)\right)
 \big/\left\langle
 [\lambda\circ\omega(f)\otimes v]_M
 -[\lambda\otimes\omega(f)v]_N\right\rangle,
\]
where $f:M\to N$, $v\in\omega(M)$, and $\lambda\in\omega(N)^*$.

\begin{proposition}[Matrix-coefficient coalgebra]\label{cyiv:k3-coend}
For a basis $(e_i)$ with dual basis $(e^i)$, the formulas
\[
 \Delta[\lambda\otimes v]_M
 =\sum_i[\lambda\otimes e_i]_M\otimes[e^i\otimes v]_M,
 \qquad \epsilon[\lambda\otimes v]_M=\lambda(v)
\]
are basis-independent and make $H$ a coalgebra.
Tensor product gives its multiplication, and duals give its antipode.
Each $\omega(M)$ is a comodule by
\[
 v\longmapsto\sum_i e_i\otimes[e^i\otimes v]_M.
\]
Every selected kernel is a comodule morphism.
\end{proposition}
\begin{proof}
The identity tensor $\sum_i e_i\otimes e^i$ is basis-independent.
Inserting it twice proves coassociativity; evaluation proves both
counit identities.
The defining relations express naturality under $f$ and are preserved
by these formulas.
The strong monoidal constraints identify the coefficient of $M\otimes N$
with the product of coefficients of $M$ and $N$.
Their pentagon and unit identities give the bialgebra axioms.
The evaluation and coevaluation identities for duals give the antipode
identities.
The same relations prove that $\omega(f)$ intertwines the coactions.
\end{proof}

The linear dual of $H$ represents linear natural endomorphisms of
$\omega$.
Tensor-compatible natural automorphisms are the corresponding
multiplicative, unit-preserving families, not all elements of this dual.
An equivalence with a comodule category requires the additional exactness
and reconstruction hypotheses stated in the Tannaka theorem below.
The action algebra $A_S$ acts by selected kernels; the coend records
symmetries commuting with those kernels.

\part{The exact \texorpdfstring{$K3\times E$}{K3 x E} protected completion}

\chapter{The protected input}
The K3 elliptic genus is $2\phi_{0,1}$.
The Borcherds half determines the weight-five denominator $D_5$ and its
Borcherds Lie superalgebra $\mathfrak g_\Delta$.
Fix its positive root spaces with their parity and proper height.
The positive algebra and its root Fock space are defined from this Lie
superalgebra, independently of a compact Hall realization.

For each finite zero-potential quiver window $Q_{\Delta,H}$, form the
Ginzburg algebra $\Gamma_{\Delta,H}$.
Restriction to full subquivers gives the projections of
Proposition~\ref{cyiv:ginzburg-projection}.
Write $\Gamma_\Delta$ for this specified inverse system.
A nondegenerate continuous Calabi--Yau structure on a completed limit
is additional data.
An ordinary Hall category $\mathcal B_\Delta$ must carry its chosen
quiver coefficients and extension correspondences.
Its identification with a root presentation requires the corresponding
Hall theorem on that precise datum.

\begin{definition}[Root-datum realizations]
A protected root-datum realization records
\[
 \mathfrak P(K3\times E)=
 (\Gamma_\Delta,\mathcal B_\Delta,\mathcal H_{\mathrm{1p}},
 \mathcal F^{\mathrm{hyb}}_{\Delta,E},\widehat U_q(\mathfrak g_\Delta),
 \mathbb A_Z).
\]
Here $\mathcal H_{\mathrm{1p}}=\bigoplus_{\alpha>0}(\mathfrak g_\Delta)_\alpha$
with its parity.
The current-factorization, paired quantum-group, and Fredholm entries
retain their own coefficient categories, completions, and defining
hypotheses.
Maps among these entries are specified realization maps, not consequences
of using the same root labels.
\end{definition}

The enveloping and Chevalley constructions are functorial under the
primitive maps of Theorem~\ref{cyiv:primitive-comparison}.
The root Fock trace is computed from the actual graded root spaces.
A Fredholm realization adds a family of operators on a specified domain
and the convergence and orientation data identifying its Pfaffian with
the denominator.
These are the analytic inputs to the square-root comparison below.

\begin{definition}[Integral protected realization]
Let $\mathfrak p_X^+$ be the BPS primitive Lie algebra of a selected strong-oriented cohomological lift of the geometric $E$-equivariant Hall source.  The integral protected realization is
\[
 \mathfrak P_{\mathrm{int}}(X,\Delta)
 =\left(
   \Hcal_X^{E,\mathrm{mot}},
   \mathfrak P(K3\times E),
   \mathfrak b_{X\boxtimes\Delta},
   r_X,r_\Delta,
   I_{X,\Pi}^E
 \right).
\]
At the motivic level the first, second, and character components are defined.  A strong orientation enhancement supplies $\mathfrak p_X^+$ and hence the amalgamated primitive component.
\end{definition}

\begin{proposition}[Primitive graph comparison]
Suppose the two primitive Borels, their common Cartan action, and their
completed amalgam are supplied as above.
Their Cartan retractions induce retractions from the amalgam to each
factor.
For each specified primitive map $\varphi$ preserving Cartan action,
its graph quotient is canonically the geometric Borel.
\end{proposition}
\begin{proof}
The identity on one Borel and the Cartan projection of the other agree
on the common Cartan, so the amalgam gives the retraction.
The quotient statement is Proposition~\ref{cyiv:graph-quotient}.
\end{proof}

The motivic integration map and primitive comparison are different maps.
Strong orientation gives the coefficient coherence of a cohomological
Hall construction under its stated geometric hypotheses.
Extracting a primitive Lie algebra additionally requires the appropriate
BPS or coproduct construction.

\chapter{Cyclic open-string operations}
Let $A$ be a finite-dimensional, uncurved cyclic $A_\infty$ algebra
over $\mathbb C$, with a nondegenerate pairing of degree $-3$.
Use the cohomological suspension $s:A\to A[1]$ of degree $-1$.
The suspended operations $b_n:(A[1])^{\otimes n}\to A[1]$ have degree one.
Our sign convention is the coderivation convention $b^2=0$ on the
tensor coalgebra, with the Koszul rule for moving homogeneous inputs.
Write $\omega$ for the resulting degree-$-1$ symplectic form on $A[1]$.
Fix the Hamiltonian convention $\iota_{X_F}\omega=dF$ and
$[X_F,X_G]=X_{\{F,G\}}$.

For the universal degree-zero coordinate $x$ on $A[1]$, put
\[
 S(x)=\sum_{n\ge1}\frac1{n+1}\omega(x,b_n(x^{\otimes n})).
\]
Cyclic invariance is read with the suspended Koszul signs.
All functions are formal power series at the origin.
An equivalent unsuspended convention writes this as
$\frac12\langle a,m_1a\rangle+
\sum_{n\ge2}\frac1{n+1}\langle a,m_n(a^{\otimes n})\rangle$.

\begin{proposition}[Cyclic master equation]\label{cyiv:cyclic-master}
The action satisfies $\{S,S\}=0$.
\end{proposition}
\begin{proof}
Cyclicity makes the $n+1$ derivatives of the arity-$n+1$ term equal,
with the suspended Koszul signs.
Thus $X_S(x)=\sum_{n\ge1}b_n(x^{\otimes n})$.
The coderivation identity implies $X_S^2=0$.
The Hamiltonian identity gives $X_{\{S,S\}}=2X_S^2=0$.
Nondegeneracy of $\omega$ makes $\{S,S\}$ constant.
Uncurvedness removes $b_0$, so $S$ starts in order two and its bracket
has no constant term.
The constant is therefore zero.
\end{proof}

A finite-height Ginzburg realization requires such a cyclic model on
the selected proper objects before this scalar action is applied.
A Hall realization additionally requires a moduli problem with
critical coefficient complexes.
For an extension $A_0\to F\to B_0$, put
$L(A_0,B_0)=\det R\operatorname{Hom}(A_0,B_0)$.
Three-Calabi--Yau duality gives
\[
 R\operatorname{Hom}(B_0,A_0)
 \simeq R\operatorname{Hom}(A_0,B_0)^\vee[-3].
\]
Applying the determinant functor to the two derived Hom triangles gives
\[
 \det R\operatorname{Hom}(F,F)
 \cong \det R\operatorname{Hom}(A_0,A_0)
 \otimes\det R\operatorname{Hom}(B_0,B_0)
 \otimes L(A_0,B_0)^{\otimes2}.
\]
Indeed, the four subquotients are the two diagonal Hom complexes and
the two cross complexes.
Dualization and the odd shift cancel on their determinant lines,
so the two cross determinants coincide.
Orientation lines require square-root isomorphisms along this formula,
compatible on triple flags and with any quotient descent.
The cyclic master equation does not select these isomorphisms.

\chapter{An algebraic open--closed comparison}
A finite-dimensional ordinary Lie algebra $\mathfrak l$ gives a
precise coefficient model for an enveloping-algebra centre.
Let $U=U(\mathfrak l)$ over $\mathbb C$.
The adjoint module $U^{\mathrm{ad}}$ has action
$x\cdot u=xu-ux$.
Its Chevalley cochains are
\[
 C^n(\mathfrak l,U^{\mathrm{ad}})
 =\operatorname{Hom}_{\mathbb C}(\Lambda^n\mathfrak l,U),
\]
with differential
\[
\begin{split}
 (df)(x_1,\ldots,x_{n+1})
 &=\sum_i(-1)^{i+1}[x_i,f(x_1,\widehat x_i,\ldots,x_{n+1})]\\
 &\quad+\sum_{i<j}(-1)^{i+j}
 f([x_i,x_j],x_1,\widehat x_i,\widehat x_j,\ldots,x_{n+1}).
\end{split}
\]
The second line places the bracket first and retains the order of the
remaining arguments.

\begin{theorem}[Hochschild--Chevalley comparison]\label{cyiv:hh-ce}
Restriction followed by antisymmetrization defines a quasi-isomorphism
\[
 C^\bullet_{\mathrm{HH}}(U,U)
 \longrightarrow C^\bullet(\mathfrak l,U^{\mathrm{ad}}),
 \qquad
 \phi\longmapsto\left[(x_1,\ldots,x_n)\mapsto
 \sum_{\sigma\in S_n}\operatorname{sgn}(\sigma)
 \phi(x_{\sigma(1)},\ldots,x_{\sigma(n)})\right].
\]
This is a comparison of ordinary cochain complexes.
\end{theorem}
\begin{proof}
The Chevalley bimodule resolution has terms
$K_n=U\otimes\Lambda^n\mathfrak l\otimes U$.
Its differential has the left-minus-right multiplication terms and
the bracket terms with the signs dual to the displayed differential.
Filter $K_n$ by total PBW degree, counting each exterior generator
as degree one.
The multiplication terms preserve this degree and bracket terms lower it.
The associated graded is the Koszul resolution of the diagonal in
$\operatorname{Sym}(\mathfrak l)\otimes\operatorname{Sym}(\mathfrak l)$,
with regular sequence $x_i\otimes1-1\otimes x_i$.
That Koszul complex is exact away from degree zero.
The filtration is bounded below and exhaustive, and each element has
finite filtration degree, so lifting a leading term proves exactness
of $K_\bullet$ by induction.
Its augmentation has image $U$.
Thus $K_\bullet$ is a free bimodule resolution.

Antisymmetrization gives a map from $K_\bullet$ to the ordinary bar
resolution, lifting the identity of $U$.
Adjacent products in its bar differential pair to commutators;
$xy-yx=[x,y]$ gives exactly the bracket terms of $K_\bullet$.
The outer products give its left-minus-right terms.
A comparison between projective resolutions lifting the identity is
a homotopy equivalence.
Applying $\operatorname{Hom}_{U\otimes U^{\mathrm{op}}}(-,U)$ gives
the stated formula and quasi-isomorphism.
\end{proof}

This is the classical Cartan--Eilenberg comparison.
For a finite Lie superalgebra, the same construction uses the
super-exterior powers and permutation Koszul signs.
Odd generators produce symmetric exterior directions, so the resolution
can be unbounded; its algebraic direct sums and cochain products remain
part of the convention.
The PBW and super-Koszul proof applies with those conventions.

A physical mixed holomorphic--topological realization requires a map
from its independently defined bulk observables to this cochain object.
The theorem specifies the adjoint coefficients and the actual algebraic
comparison.
An $E_2$ refinement, a current-factorization realization, and compatibility
with root-height transitions require maps preserving their additional
operations.
A quotient of Lie algebras alone does not provide a covariant map of
Hochschild centres.

\chapter{The exact scalar and first-order observables}
Set
\[
 \DenIg(Z)=D_5(2Z),
 \qquad \jmath(Z)=Z/2,
 \qquad \jmath^*\DenIg=D_5.
\]
The one-particle root spaces give
\[
 \sch\Fock(\gDelta_+)
 =\prod_{\alpha>0}(1-e^{-\alpha})^{-\sdim\gDelta_\alpha}.
\]
The oriented root operator gives
\[
 e^{-\rho}\operatorname{sPf}_F(\mathbb A_Z)=D_5(2Z).
\]
The primitive reduced curve-counting theorem gives
\[
 \mathcal Z^{\red}_{K3\times E}=-(\jmath^*\DenIg)^{-2}=-\chi_{10}^{-1}.
\]

\begin{theorem}[Protected square-root theorem]
Assume the stated root-space Fredholm comparison with its orientation,
Weyl shift, and convergence domain.
Within that realization, the denominator section $\DenIg$ is simultaneously:
\begin{enumerate}
\item the Weyl denominator of $\gDelta$;
\item the Euler character of the bar--Chevalley complex;
\item the oriented Fredholm Pfaffian of $\mathbb A_Z$;
\item the PBW determinant of the parity-refined root superspace $\mathcal H_{\mathrm{1p}}$;
\item the automorphic denominator whose half-coordinate pullback squares to the reduced scalar denominator.
\end{enumerate}
All five identifications preserve root degree and Weyl transport.
\end{theorem}
\begin{proof}
The first two are the denominator and bar--Chevalley theorems.  The third is the Fredholm theorem.  PBW for the Borcherds Lie superalgebra identifies the root superspace determinant, giving the fourth.  The half-coordinate identity $\jmath^*\DenIg=D_5$, the theta normalization, and the reduced curve-counting theorem give the fifth.
\end{proof}

\chapter{Physical meaning}
The root-datum realizations retain several distinct constructions.
The finite Ginzburg algebras and their projections form an inverse system.
A completed three-Calabi--Yau sector requires a continuous nondegenerate
Calabi--Yau structure and compatible cyclic models on the selected proper objects.
The quiver Hall category retains its chosen coefficients and extension
correspondences. A Hall theorem on that datum must identify its root presentation.
The Borcherds root superspace defines the root-graded Fock space and its
supercharacter. Its identification with physical one-particle states
requires a parity-preserving graded comparison intertwining the primitive actions.

A realization on the elliptic curve must construct the current object
with its local and wrapped supports, collision operations, and comparison maps.
The quantum Borcherds algebra, with its specified pairing and completion,
acts as a line-operator symmetry only through a representation preserving
fusion and exchange. The root operator gives the Fredholm product on its
specified convergence domain. A physical determinant interpretation requires
an operator comparison and an orientation-preserving identification of the
determinant lines and sections. The Fock character gives the scalar square.
An operator trace additionally requires its state space and trace domain.

Simple-root states satisfying the defining relations give a primitive Lie map.
The explicit enveloping and Chevalley extensions preserve products and brackets.
A compact string realization must additionally construct its brane functor,
Hall orientation maps, paired quantum coproduct, and trace-line comparison.

\chapter{Primitive comparisons and their algebraic extensions}
Let $\mathfrak n_\Delta$ be the positive Borcherds Lie superalgebra.
Choose homogeneous parity-preserving images of its simple generators
in a charge-graded Lie superalgebra $\mathfrak p$.
Assume these images satisfy every defining relation.

\begin{theorem}[Extension of a primitive comparison]\label{cyiv:primitive-comparison}
The chosen images determine a unique Lie-superalgebra map
$f:\mathfrak n_\Delta\to\mathfrak p$.
It induces
\[
 U(f)(x_1\cdots x_r)=f(x_1)\cdots f(x_r),
 \qquad C_*(f)=\operatorname{Sym}^c(f[1]).
\]
These are respectively an enveloping-algebra map and a Chevalley
chain-coalgebra map.
Filtration-preserving maps extend to the specified inverse-limit
completions.
If $f$ is surjective and both algebras have finite charge components
and proper positive height, equality of their parity-refined PBW
characters implies that $f$ and $U(f)$ are isomorphisms.
\end{theorem}
\begin{proof}
The defining presentation gives $f$.
Its tensor-algebra extension carries
$xy-(-1)^{|x||y|}yx-[x,y]$ to the corresponding relation.
It therefore descends to $U(f)$.
The shifted symmetric-coalgebra map commutes with the Chevalley
differential because $f$ preserves brackets.
Compatible maps on filtration quotients give maps on their inverse limits.
Surjectivity of $f$ gives surjectivity of $U(f)$.
Character equality makes each finite charge component an isomorphism,
separately in both parities.
The PBW embeddings then imply injectivity of $f$.
\end{proof}

At level zero the current Lie conformal algebra is
$\mathbb C[\partial]\otimes\mathfrak n_\Delta$ with
$[x_\lambda y]=[x,y]$ on constant generators.
The map $\partial^j x\mapsto\partial^j f(x)$ preserves this bracket
and sesquilinearity.
A nonzero central level additionally requires preservation of the
specified invariant pairing.
Hall convolution, Calabi--Yau structures, quantum coproducts, and
determinant lines require their own comparison data.

\chapter{Finite quivers and Calabi--Yau completions}
For a finite quiver $Q$ with zero potential, let $\Gamma(Q,0)$ be its
path-length-completed Ginzburg algebra.
Path products are read from left to right: if $a:i\to j$ and
$b:j\to k$, then $ab:i\to k$.
For each arrow $a:i\to j$, add $a^*:j\to i$ of degree $-1$.
At each vertex add a loop $t_i$ of degree $-2$.
Original arrows have degree zero, and
\[
 da=da^*=0,\qquad
 dt_i=e_i\left(\sum_a[a,a^*]\right)e_i.
\]

\begin{proposition}[Restriction to a full subquiver]\label{cyiv:ginzburg-projection}
If $Q_0$ is a full subquiver, retaining its generators and killing all
other vertices, arrows, reverse arrows, and loops defines a continuous
dg-algebra projection
\[
 \Gamma(Q,0)\longrightarrow\Gamma(Q_0,0).
\]
The inclusion of retained graded generators need not commute with
the differential.
\end{proposition}
\begin{proof}
The projection respects source and target idempotents and extends
multiplicatively.
At a retained vertex, every term involving a discarded arrow vanishes.
The other terms are precisely $dt_i$ in the smaller algebra.
All arrow differentials vanish, so the projection is a chain map.
It preserves the path-length filtration and is continuous.
For the final assertion, start with one vertex $i$ and no arrows,
where $dt_i=0$.
Add a vertex $j$ and an arrow $a:i\to j$.
In the larger algebra $dt_i=aa^*$, a nonzero length-two path.
Thus the graded inclusion is not a dg map.
\end{proof}

Ginzburg completion supplies a three-Calabi--Yau algebra for each finite
quiver in its stated setting.
The displayed projections specify an inverse system for compatible
zero-potential windows.
A Calabi--Yau structure on a completed limit requires its own
nondegeneracy and continuity argument.
Comparisons with Hall categories and quantized Borcherds algebras must
also preserve orientations, root relations, and the paired Hopf data.

\part{Quantum groups, exchange, and representation theory}

\chapter{From positive half to full quantum group}
Suppose the compact or local positive algebra $H^+$ is paired continuously with a negative half $H^-$, carries a completed coproduct, and admits Cartan data and a separated triangular completion. The completed double
\[
 \qgroup=\widehat D(H^+)
\]
has triangular decomposition $H^-\widehattensor H^0\widehattensor H^+$. Its universal $R$-matrix acts on continuous tensor products of modules. The classical limit is a Lie bialgebra whose cobracket is the first-order term of the coproduct asymmetry.

For $\C^3$, $H^+$ is the positive affine Yangian. For local ADE surfaces, it is the completed positive affine ADE Yangian. For the Igusa algebra, the paired Borel construction of the preceding part supplies a height-complete formal quantum group; a compact Hall realization maps to it through primitive, coproduct, and pairing comparisons.

\chapter{Tannaka reconstruction}
Let $\mathcal L$ be an essentially small rigid $\kfield$-linear abelian monoidal category with finite-dimensional morphism spaces, finite-length objects, and exact tensor product. Let
\[
 \omega:\mathcal L\to\Vect_{\mathrm{fd}}
\]
be a faithful exact strong monoidal functor. Assume the coend
\[
 H=\int^{M\in\mathcal L}\omega(M)^*\tensor\omega(M)
\]
exists and that the canonical comparison functor to finite-dimensional $H$-comodules is comonadic. Tensor product makes $H$ a bialgebra, and rigidity supplies its antipode. A braiding gives a coquasitriangular form; a meromorphic braiding gives its spectral version.

\begin{theorem}[Rigid Tannaka reconstruction]
Under the preceding finite and comonadic hypotheses, the comparison functor is a monoidal equivalence
\[
 \mathcal L\simeq\operatorname{Comod}_{\mathrm{fd}}(H).
\]
Braiding transports to the universal coquasitriangular form on $H$.
\end{theorem}
\begin{proof}
Matrix coefficients define the universal $H$-coaction on every $\omega(M)$. Tensor product and duality give the bialgebra and antipode structures. The Barr--Beck comonadic theorem identifies $\mathcal L$ with coalgebras for the induced comonad, and the coend identifies that comonad with $H\tensor-$. The hexagon identities transport to the coquasitriangular identities.
\end{proof}

This reconstruction explains why a line-operator category, a fiber functor, and exchange together produce a quantum group. The ordered boundary algebra provides ordered composition.  A compatible bialgebra coproduct supplies tensor products of its modules, and exchange supplies additional braiding data.

\chapter{KZ and rational exchange}
For a simple Lie algebra $\mathfrak g$, the Knizhnik--Zamolodchikov connection is
\[
 \nabla=d-\frac{\hbar}{2\pi i}
 \sum_{i<j}\Omega_{ij}\,d\log(z_i-z_j).
\]
Flatness follows from the infinitesimal braid relations. Drinfeld--Kohno identifies its monodromy tensor category with the corresponding Drinfeld--Jimbo category after a choice of associator.

Rational difference exchange instead gives the Yangian $R$-matrix
\[
 R(u)=1-\frac{\hbar P}{u}
\]
in the fundamental representation. The Yang--Baxter equation is the confluence identity for three exchanges. The KZ and Yangian constructions share the geometry of configuration spaces while using different flat connections.

\chapter{The centre and the double}
The Hochschild centre $Z_{\ord}(A)$ controls central operations on the
ordered algebra in its specified coefficient category.
To form a Drinfeld centre of modules, first specify a monoidal module category.
For an ordered chiral bialgebra satisfying
Definition~\ref{def:vol4-ordered-chiral-bialgebra}, use
\[
 \mathcal M_A=\operatorname{Mod}_A(\widehat{\mathcal V}_C)
\]
with the monoidal product of Theorem~\ref{thm:vol4-representation-fusion}.
An object of $Z(\mathcal M_A)$ is a module $M$ together with natural equivalences
\[
 c_{M,N}:M\widehat{\tensor}^{!}N
 \xrightarrow{\sim}N\widehat{\tensor}^{!}M
 \qquad(N\in\mathcal M_A)
\]
satisfying the half-braiding hexagon and unit identities, with their coherences.
This categorical centre is braided.
A comparison with the Hochschild centre must specify a realization functor,
the source and target of the comparison, and the operations it preserves.
Dualizability alone does not supply such a comparison.
Derived loop-space descriptions apply in the geometric categories and
under the hypotheses of their respective centre theorems
\cite{BenZviFrancisNadler}.

The centre and the double answer complementary questions.
The centre classifies central actions in the chosen category.
The double reconstructs a quantum group from paired positive and negative halves.
A field-theory comparison requires explicit realizations of the relevant
line and bulk operations and a map compatible with both constructions.

\part{A compact six-dimensional protected source}

\chapter{Holomorphic cotangent BF theory on \texorpdfstring{$K3\times E$}{K3 x E}}
Let $S$ be a K3 surface with holomorphic symplectic form $\sigma_S$, let $E$ be an elliptic curve with holomorphic form $\dd z$, and put
\[
 X=S\times E,
 \qquad
 \Omega_X=\sigma_S\wedge\dd z.
\]
For a finite proper-height quotient $\nIg^{(H)}$ of the positive Igusa Borcherds algebra, take fields
\[
 A\in\Omega^{0,\bullet}(X,\nIg^{(H)})[1],
 \qquad
 B\in\Omega^{0,\bullet}(X,(\nIg^{(H)})^\vee)[1].
\]
The invariant evaluation pairing and $\Omega_X$ define the BV action
\begin{equation}\label{eq:vol4-compact-hbf}
 S_H(A,B)=
 \int_X\Omega_X\wedge
 \angles{B,\dbar A+\frac12[A,A]}.
\end{equation}
This is the holomorphic cotangent theory of the dg Lie algebra of $\nIg^{(H)}$-valued Dolbeault forms.

\begin{theorem}[Classical compact-background source]
The action \eqref{eq:vol4-compact-hbf} satisfies the classical master equation and defines a holomorphic factorization algebra on the compact Calabi--Yau threefold $X$.
\end{theorem}
\begin{proof}
The shifted cotangent pairing is invariant.  The classical master equation is the Jacobi identity of $\nIg^{(H)}$ together with $\dbar^2=0$ and the derivation property of $\dbar$.  Local observables of an elliptic cyclic $L_\infty$ algebra form a factorization algebra by the standard BV construction.
\end{proof}

\chapter{Tree exactness and the quantum master equation}
Every interaction vertex in \eqref{eq:vol4-compact-hbf} has one $B$-leg and two $A$-legs.  The propagator pairs an $A$-leg with a $B$-leg.

\begin{lemma}[Loop bound]
A connected Feynman graph of the cotangent theory has loop number at most one.
\end{lemma}
\begin{proof}
A graph with $V$ interaction vertices has exactly $V$ internal or external $B$-half-edges.  Every internal edge consumes one of them, hence $E\le V$.  The loop number is $E-V+1\le1$.
\end{proof}

The positive algebra is graded by strictly positive root height.  For homogeneous $x\in\nIg^{(H)}$, the operator $\ad_x$ raises height.  Its coadjoint operator lowers dual height.  A product of adjoint operators labelled by positive roots has total nonzero height shift and therefore has zero trace.

\begin{proposition}[Vanishing of one-loop Lie weights]\label{cyiv:loop-weights}
For every finite positive-height quotient, the Lie weight of each
one-loop graph in the cotangent theory vanishes.
\end{proposition}
\begin{proof}
A one-loop graph has $E=V$.
Every $B$-leg is consumed by an internal edge, so its external legs are
$A$-legs.
Its Lie weight is a cyclic trace of adjoint operators, or the dual
trace with the conventionally equivalent signs.
Each external label has positive root height.
The product therefore strictly raises height on the finite-dimensional
adjoint representation and has zero trace.
The same argument applies to the dual representation.
\end{proof}

This vanishing removes the one-loop Lie coefficients in a perturbative
construction with the stated propagator and vertex types.
A renormalized quantum factorization algebra still requires admissible
analytic kernels, local extensions, the quantum master equation, and
compatible height transitions on the chosen observable spaces.
These are the remaining quantum realization data.

\chapter{Classical reduction along K3}
Choose a K\"ahler metric on $S$ and normalize a closed harmonic
$(0,2)$-form $\eta$ by $\int_S\sigma_S\wedge\eta=1$.
Let $K=\mathbb C[\eta]/(\eta^2)$ with $|\eta|=2$ and zero differential.
There is a unital dg-algebra map
\[
 i:K\longrightarrow\Omega^{0,*}(S),\qquad \eta\longmapsto\eta.
\]
Its image is closed under multiplication because $(0,4)$-forms vanish.
The Dolbeault cohomology of a K3 surface is one-dimensional in degrees
zero and two and zero in degree one, so $i$ is a quasi-isomorphism.
It preserves the trace defined by integration against $\sigma_S$.

Let $\mathfrak l$ be a finite-dimensional Lie algebra, and write
$\mathfrak d=\mathfrak l\ltimes\mathfrak l^*$ for its cotangent Lie
algebra, with invariant evaluation pairing.
For an open $U\subset E$, tensor with $\Omega^{0,*}(U)$, using the
completed smooth tensor product in the $S$ direction.
The field comparison is
\[
 K\otimes\Omega^{0,*}(U)\otimes\mathfrak d
 \xrightarrow{i\otimes\id}
 \Omega^{0,*}(S\times U)\otimes\mathfrak d.
\]
The bracket is multiplication of forms followed by the Lie bracket,
with Koszul signs.

\begin{proposition}[Classical field comparison]\label{cyiv:k3-classical-reduction}
The displayed map is a trace-preserving dg Lie quasi-isomorphism,
compatible with restriction in $U$.
Its induced minimal operations in the K3 factor have no higher products.
\end{proposition}
\begin{proof}
Multiplicativity of $i$ proves bracket compatibility.
Hodge projection $p$ and $h=\bar\partial^*G$ on the compact surface satisfy
$pi=\id$ and $\bar\partial h+h\bar\partial=\id-ip$.
Tensoring these continuous operators with forms on $U$ gives the same
contraction, with the tensor-product signs.
Thus $i\otimes\id$ is a quasi-isomorphism.
Restriction acts only on the $U$ factor and commutes with it.
The normalized trace on $K$ is the restriction of the Dolbeault trace.
Choose the contraction with $hi=0$.
In every higher transferred tree, the first product of inputs from
$i(K)$ remains in $i(K)$ and the next internal homotopy kills it.
Thus every higher transferred product vanishes.
\end{proof}

The coefficient algebra $K$ retains both the degree-zero and degree-two
K3 modes.
The proposition compares classical field complexes and their brackets.
Constructing an equivalence of quantum observable factorization algebras
requires compatible quantization, support conventions, and the induced
maps on observable completions.
The ordinary Hochschild comparison of Theorem~\ref{cyiv:hh-ce} is a
separate algebraic calculation, not that quantum field map.

\chapter{Protected states and the Igusa trace}
Let
\[
 \mathcal H_{\mathrm{1p}}
 =\widehat\bigoplus_{\alpha>0}\gIg_\alpha
\]
with its Borcherds parity.  The multiparticle state space is the completed super-Fock space
\[
 \mathcal H_{\mathrm{Fock}}=\widehat\Sym(\mathcal H_{\mathrm{1p},\bar0})
 \tensor\widehat\Lambda(\mathcal H_{\mathrm{1p},\bar1}).
\]
For $Z$ in the absolute-convergence chamber, let $e^{-H_Z}$ act by $e^{-\alpha(Z)}$ on the root space of degree $\alpha$.

\begin{theorem}[Protected Hamiltonian trace]
The supertrace is
\[
 \str_{\mathcal H_{\mathrm{Fock}}}(e^{-H_Z})
 =\prod_{\alpha>0}
 (1-e^{-\alpha(Z)})^{-\sdim\gIg_\alpha}.
\]
After multiplication by the inverse Weyl line, its reciprocal is $\DenIg(Z)=D_5(2Z)$.  After the half-coordinate pullback and doubling, the disconnected protected scalar is $-(\jmath^*\DenIg)^{-2}=-\chi_{10}^{-1}$ in the Oberdieck--Pixton normalization.
\end{theorem}
\begin{proof}
An even oscillator contributes $(1-e^{-\alpha})^{-1}$ and an odd oscillator contributes $1-e^{-\alpha}$.  Multiplication over the root basis gives the product.  The denominator identity and the reduced curve-counting theorem supply the final two equalities.
\end{proof}

\chapter{The compact-source comparison}
The compact threefold provides the classical cotangent field theory
and the dg Lie comparison of Proposition~\ref{cyiv:k3-classical-reduction}.
The root datum separately gives the Borcherds algebra, its enveloping
algebra, and the root-graded Fock space.
The latter has the trace computed above.
Proposition~\ref{cyiv:loop-weights} supplies the finite one-loop Lie
cancellation, while Theorem~\ref{cyiv:hh-ce} computes the ordinary
algebraic centre model with adjoint coefficients.

\begin{problem}[Geometric compact comparison]\label{cyiv:compact-comparison}
Construct objects $E_i$ representing the simple branes and maps on their
derived morphism complexes that preserve units, composition, and higher
composition homotopies.
Realize these maps on extension stacks with the reduced critical
coefficients and orientation isomorphisms on triple flags.
Construct compatible tensor and unit constraints on the resulting
boundary functor.
Compare the chosen quantum bulk action with the independently defined
boundary centre, including support, collisions, and completion.
Finally construct a trace and a map of its value line to the automorphic
line in the same charge normalization.
\end{problem}

A primitive map satisfying the root relations gives the enveloping and
Chevalley maps of Theorem~\ref{cyiv:primitive-comparison}.
A strong monoidal equivalence transports categorical centres.
Neither construction determines the brane functor or its cyclic trace.
These are the additional maps required for compact geometric
reconstruction.

\part[Physics]{Physics of the comparison maps}

\chapter{Holomorphic--topological boundary theories}
A three-dimensional holomorphic--topological theory on $C\times\R$ has local observables that are holomorphic along $C$ and locally constant along $\R$. A boundary produces an ordered chiral algebra. A second topological direction permits exchange and gives an $\Etwo$ enhancement or a meromorphic braided line category. The bulk acts through the derived centre.

The finite cotangent BF model gives the exact local calculation:
\[
 \Barc(U(\mathfrak l))\simeq C_*(\mathfrak l).
\]
For $\mathfrak l=\nDelta^{(H)}$, this calculation realizes every finite-height Igusa denominator sector.

\chapter{Six-dimensional hCS on \texorpdfstring{$K3\times E$}{K3 x E}}
Choose a holomorphic volume form $\Omega_S\wedge dz_E$. The hCS theory has classical fields
\[
 \Omega^{0,*}(S\times E,\mathfrak a)[1]
\]
and action $S_{\hCS}$. Its factorization observables are native to the threefold. Pushforward along $p_E$ gives a factorization algebra on $E$ whose coefficient complex contains the Dolbeault modes of $S$.

A protected reduction to a chiral algebra consists of a deformation retract of the K3 Dolbeault complex, compatibility of the transferred $L_\infty$ operations with the trace, a boundary or defect selecting ordered fusion, and a quantization trivializing the local anomaly. These data produce the ordered chiral algebra to which the Hall/BPS comparison may be applied.

The Hodge equality $\chi(\Ocal_{S\times E})=0$ controls one compact supertrace. The gauge and gravitational anomalies are classes in the local BV deformation complex and are evaluated by their own characteristic forms.

\chapter{Omega backgrounds}
Let a torus $T$ act holomorphically on the local geometry and preserve the selected supercharge. For a generator $V$ of the action, the equivariant differential is
\[
 Q_\Omega=Q+\iota_V,
 \qquad Q_\Omega^2=\mathcal L_V.
\]
After localization, fixed loci produce deformation parameters and often a spectral $R$-matrix. Toric threefolds and ADE surface resolutions carry such actions. A deformation family equipped with monodromy-compatible moduli correspondences and wall-crossing equivalences transports the localized algebra from algebraic or equivariant K3 fibers to the selected generic fiber.

\chapter{Anomalies as typed classes}
The principal anomaly classes are
\[
\begin{array}{c|c}
\text{construction}&\text{deformation complex}\\ \hline
\text{hCS quantization}&\text{local BV cochains}\\
\text{critical CoHA orientation}&\text{Picard/gerbe theory of the critical stack}\\
\text{Hall coproduct}&\text{correspondence and localization complex}\\
\text{projective conformal blocks}&\text{Atiyah algebra of the determinant line}\\
\text{automorphic multiplier}&\text{character group of the modular cover}.
\end{array}
\]
A comparison theorem is a morphism between the corresponding complexes together with equality of the transported classes. This formulation turns numerical coincidences into geometric identities precisely when an intertwining map exists.

\part[Automorphic forms]{Automorphic characters and scalar shadows}

\chapter{Borcherds products as characters}
Let $\mathfrak n_+$ be a locally finite positive Lie superalgebra. The Chevalley Euler character is
\[
 \chi C_*(\mathfrak n_+)
 =\prod_{\alpha>0}(1-e^{-\alpha})^{\sdim\mathfrak g_\alpha}.
\]
A Borcherds product becomes the Weyl-shifted character of $C_*(\mathfrak n_+)$. This construction produces a character of an existing algebra. Hall geometry supplies a second algebra, and a primitive comparison theorem identifies the two.

\chapter{The universal Borcherds weight formula}
For a weight-zero weak Jacobi form
\[
 \varphi(\tau,z)=\sum c(n,\ell)q^nr^\ell,
\]
the Borcherds lift has weight
\[
 \wt(\operatorname{Borch}(\varphi))=\frac{c(0,0)}2.
\]
For the K3 half-index, $c(0,0)=10$ and the weight is $5$. This formula lives entirely on the automorphic layer. Hodge Euler characteristics, Heisenberg ranks, central charges, and BV anomalies are evaluated by their own traces.

\chapter{The Gritsenko--Cl\'ery atlas}
The eight diagonal-divisor forms are
\[
 \Delta_5,\Delta_2,\Delta_1,\Delta_{1/2},
 \nabla_3,\nabla_2,\nabla_{3/2},Q_1
\]
with weights
\[
 5,2,1,\frac12,3,2,\frac32,1.
\]
Each form is an automorphic product on its own paramodular or Hecke congruence group with its own character or multiplier \cite{GC}. A geometric realization is specified row by row by a charge lattice, Hall problem, trace, and primitive algebra comparison.

\chapter{Scalarization hierarchy}
The constructions descend through four levels:
\[
\begin{aligned}
 \text{factorization or Hall object}
 &\longrightarrow \text{primitive Lie algebra}\\
 &\longrightarrow \text{PBW character}
 \longrightarrow \text{scalar partition function}.
\end{aligned}
\]
Each arrow retains the quotient of structure displayed at its target. The first object contains supports, extension maps, and brackets. The second retains the infinitesimal algebra. The third retains signed graded dimensions. The fourth adds a trace and normalization. Upward reconstruction is supplied by the corresponding geometric operations and comparison morphisms.

\part{The unified theorem system}

\chapter{The foundational theorem}
\begin{definition}[Complete Calabi--Yau quantum-group datum]
A complete comparison datum specifies independent categorical, field,
boundary, Hall, Hopf, and trace objects together with:
\begin{enumerate}
\item PTVV-representable moduli for the smooth proper Calabi--Yau category;
\item a renormalized BV quantization satisfying the quantum master equation and a local compact factorization generator $b$;
\item a boundary reduction or proper locally trivial map to a curve satisfying factorization pushforward, together with the supplied curve-level faces and aligned distributive law;
\item a strong-oriented Hall coefficient theory whose pull--push maps satisfy base change and whose positive algebra is separated and complete;
\item paired negative and Cartan halves, a completed coproduct, a continuous nondegenerate Hopf pairing, and a separated triangular completion on which the canonical tensors converge;
\item compatible open--closed and Hall--boundary morphisms;
\item a parity-refined primitive comparison for every automorphic identification.
\end{enumerate}
\end{definition}

\begin{theorem}[Calabi--Yau quantum-group assembly]
Let a complete datum consist of:
\begin{enumerate}
\item a Calabi--Yau category and a representable moduli-of-objects problem;
\item a quantized cyclic elliptic BV theory and a boundary or defect generator;
\item a Hall extension correspondence with coefficient functor, base-change theorem, strong orientation data, and separated charge completion;
\item paired positive and negative bialgebras, Cartan data, a nondegenerate continuous Hopf pairing, and a completed double;
\item every curve-level distributive, PBW, and analytic comparison used below.
\end{enumerate}
Then the datum determines functorially:
\begin{enumerate}
\item framed $\Etwo$ Hochschild cochains and the shifted-symplectic moduli structure;
\item the three curve-level algebraic faces and their three bars;
\item the positive Hall/BPS algebra and its primitive Lie algebra;
\item the completed quantum double and braided representation category;
\item every supplied comparison among categorical, field-theoretic, Hall, current, quantum, and automorphic realizations;
\item the Weyl-shifted PBW/Chevalley character of any parity-refined primitive Lie-superalgebra identification.
\end{enumerate}
\end{theorem}
\begin{proof}
Cyclic Deligne and shifted symplectic geometry give the first output \cite{BravRozenblyum,PTVV,BravDyckerhoff}.  Boundary factorization gives the curve faces, and Volume~I gives their bars and centres.  The Hall correspondence and its base-change theorem give the associative product; strong orientation supplies the coefficient coherence required in the $3$-Calabi--Yau setting \cite{KPS}.  The paired bialgebra data give the double and universal $R$-matrix.  The comparison maps commute by hypothesis.  PBW and Chevalley antisymmetrization give the final character.  Each construction is natural under a morphism preserving the displayed data.
\end{proof}

The local quiver, equivariant surface, K3 correspondence, and Igusa Borcherds constructions have the carriers stated above.
A geometric $K3\times E$ comparison requires the primitive map and determinant-line map of Problem~\ref{cyiv:compact-comparison}.

\chapter{Landscape of established realizations}
\begin{center}
\footnotesize
\begin{tabularx}{0.99\textwidth}{@{}>{\raggedright\arraybackslash}p{0.15\textwidth}>{\raggedright\arraybackslash}p{0.25\textwidth}>{\raggedright\arraybackslash}X>{\raggedright\arraybackslash}p{0.20\textwidth}@{}}
\toprule
\textbf{Geometry}&\textbf{Positive algebra}&\textbf{Full or centred object}&\textbf{Character}\\
\midrule
$A_2$ quiver&$U_v^+(\mathfrak{sl}_3)$&$U_v(\mathfrak{sl}_3)$&PBW root product\\
$\C^3$&$Y^+(\widehat{\mathfrak{gl}}_1)$&$Y(\widehat{\mathfrak{gl}}_1)$&plane-partition and Fock traces\\
$\widetilde{\C^2/\Gamma}$&$Y^+(\widehat{\mathfrak g}_\Gamma)$&paired completion when the negative half and coproduct are supplied&surface Fock traces\\
$K3$&$U(\mathfrak g_{\BPS}^+)$&Heisenberg and LLV actions&K3 Hall characters\\
$K3\times E$&Borcherds root envelope; geometric reduced character&$U_q(\gDelta)$ and the root-current model and classical BF theory&$-\ChiTen^{-1}$\\
\bottomrule
\end{tabularx}
\end{center}
The $A_2$ and $\C^3$ rows provide complete standard quantum-group models. The ADE row provides an established completed positive half together with its geometric current actions; its paired full double is formed from the negative half, Cartan data, coproduct, continuous Hopf pairing, and triangular completion stated in the table. The K3 row provides global two-Calabi--Yau BPS and cohomological symmetries. The final row records independently constructed root, character, and classical field objects.
Their geometric Hall realization requires the comparison in Problem~\ref{cyiv:compact-comparison}.

\chapter{Constructed comparisons and their geometric extensions}
The Hall flag argument constructs associativity once its coefficient
and orientation maps are defined.
The tensor-action formula constructs representation fusion from a
coassociative counital coproduct.
The matrix-coefficient formulas construct the coend of a supplied
fibre functor.
The primitive presentation gives enveloping and Chevalley maps.
The classical cyclic and K3 field calculations retain their actual
differentials, traces, and brackets.

The quantum field realization requires admissible quantization and
observable completions.
The compact Hall realization requires the brane and orientation maps
of Problem~\ref{cyiv:compact-comparison}.
The Fredholm and automorphic comparisons require their analytic domains
and value-line isomorphisms.
These constructions preserve the full geometric reconstruction problem
beyond its finite algebraic and classical calculations.

\chapter{Actions and comparison maps}
For a specified boundary algebra $A_b$, a coherent bulk action is a map
\[
 \beta:\Obs_{\mathrm{bulk}}|_\partial\longrightarrow Z_{\ord}(A_b).
\]
The target is the derived endomorphism object of the diagonal bimodule
in the chosen factorization coefficient category.
A Hall construction instead starts with
\[
 \M_\alpha\times\M_\beta\xleftarrow{p}\mathcal E_{\alpha,\beta}
 \xrightarrow{q}\M_{\alpha+\beta}.
\]
Critical coefficients, orientations, and admissible pull--push maps
give its product.
Framed moduli supply a representation after their action
correspondences and triple-flag identities have been constructed.
A paired Hopf structure gives $D(H^+,H^0,H^-)$.
Its coproduct defines tensor products of continuous modules, and a
convergent universal $R$-matrix gives exchange.

A strong monoidal equivalence between this module category and a
boundary category transports categorical centres.
It does not identify that categorical centre with the object
$Z_{\ord}(A_b)$ without a realization theorem.
A trace on a specified representation then gives an amplitude.
Its comparison with an automorphic section requires a map of value
lines with the same normalization.

\chapter{Comparison data for Calabi--Yau constructions}
Comparisons among Calabi--Yau categories, holomorphic field theory, chiral algebras, Hall algebras, quantum groups, and automorphic products require the following structures:
\begin{enumerate}
\item for a smooth proper $d$-Calabi--Yau category with its nondegenerate negative-cyclic class, the categorical outputs are framed $\Etwo$ Hochschild cochains and, when the moduli stack of objects is representable in the PTVV setting, $(2-d)$-shifted symplectic moduli; the shift $d$ does not change the Hochschild operadic level to $\En$;
\item a geometric BV theory and boundary category produce factorization observables and transverse ordered chiral algebras; associative chiral convolution is recorded as a separate supplied structure;
\item pushforward to a curve uses an actual map $p:X\to C$; formal cycle-specialization notation alone specifies no pushforward;
\item extension stacks, coefficient sheaves, orientations, and completions define the Hall algebra;
\item coproduct, Cartan, Hopf pairing, and topology define the double;
\item K3 Heisenberg, LLV orthogonal symmetry, local ADE Yangians, and global BPS generalized Kac--Moody algebras occupy distinct strata joined by explicit comparison maps where constructed;
\item the graph of the $K3\times E$ Gram map becomes multiplicative in the polarized central extension;
\item point-local and wrapped sectors form a hybrid factorization problem;
\item the Igusa form is the Weyl-shifted PBW/Chevalley character of $\gDelta$, and compact Hall geometry maps to it through the primitive comparison;
\item anomaly identities are transported through explicit maps between deformation complexes.
\end{enumerate}
The categorical, geometric, Hall, and automorphic operations thus have separate carriers and comparison maps.

\part{Technical appendices to the preceding book}
\chapter{Typed data table}
\begin{center}
\small
\begin{tabularx}{0.98\textwidth}{@{}>{\raggedright\arraybackslash}p{0.20\textwidth}>{\raggedright\arraybackslash}X>{\raggedright\arraybackslash}p{0.28\textwidth}@{}}
\toprule
\textbf{Object}&\textbf{Minimal data}&\textbf{Primary operation}\\
\midrule
$\CY$ category&nondegenerate negative-cyclic trace&Hochschild and shifted-symplectic structures\\
Factorization algebra&local observables and Weiss descent&disjoint-open product\\
Transversely ordered chiral algebra&factorization coefficient with an internal $\Eone$ structure for $\tensor^!$&ordered fusion\\
Associative chiral algebra&an $\Eone$ structure for chiral convolution $\starCh$&nonlocal chiral composition\\
$\CoHA$&extension correspondence and coefficient theory&pull--push product\\
Bialgebra&compatible product and coassociative counital coproduct&fusion of representations\\
Double&positive and negative halves, Cartan data, Hopf pairing&universal $R$-matrix\\
Borcherds product&Jacobi or vector-valued modular input&graded character\\
\bottomrule
\end{tabularx}
\end{center}

\chapter{The three K3 brackets}
The Heisenberg bracket is the excess intersection of two Nakajima incidence kernels.  The LLV bracket is the commutator of diagonal Lefschetz kernels and their adjoints.  The local Yangian bracket is the convolution commutator of ADE Hecke kernels with spectral equivariance.  In $\operatorname{Corr}_{K3}(S)$ all three are compositions of geometric kernels.  Their mixed commutators are the corresponding triple-intersection classes.  The action algebra and Tannaka coend retain these classes while preserving the three geometric constructions.

\chapter{Exact arithmetic checks}
The companion verification calculation establishes:
\begin{enumerate}
\item the low Fourier rows of $\phi_{0,1}$ through $q^3$ by exact theta-series division;
\item the vector-valued cocycle identity for the Gram polarization;
\item the rational Yang--Baxter equation on $(\C^2)^{\tensor3}$;
\item the $A_2$ Hall flag counts at ranks zero and one;
\item the type-II Gram matrix and reflection identities.
\end{enumerate}
Each calculation uses integer or rational arithmetic and appears as a worked formula in the relevant chapter.
